\documentclass[11pt,a4paper]{article}
\usepackage[T1]{fontenc}
\usepackage[utf8]{inputenc}
\usepackage{newtxtext}
\usepackage{amsmath,amsthm,mathtools}
\usepackage{newtxmath}
\usepackage[margin=27mm,headheight=15pt]{geometry}
\usepackage{microtype}
\usepackage{booktabs,tabularx,array}
\usepackage[dvipsnames]{xcolor}
\usepackage{enumitem,etoolbox,fancyhdr}
\usepackage{hyperref,bookmark}
\hypersetup{colorlinks=true,linkcolor=MidnightBlue,citecolor=MidnightBlue,
 urlcolor=MidnightBlue,pdftitle={Affine Monodromy and Spherical Principal Parts in Quaternionic Fueter Inversion},
 pdfauthor={Research manuscript prepared with ChatGPT},
 pdfsubject={Global inverse Fueter theorem, periods, holomorphic encoding, spherical singularities}}
\urlstyle{same}
\numberwithin{equation}{section}
\newtheorem{theorem}{Theorem}[section]
\newtheorem{proposition}[theorem]{Proposition}
\newtheorem{lemma}[theorem]{Lemma}
\newtheorem{corollary}[theorem]{Corollary}
\theoremstyle{definition}
\newtheorem{definition}[theorem]{Definition}
\newtheorem{example}[theorem]{Example}
\theoremstyle{remark}
\newtheorem{remark}[theorem]{Remark}
\newcommand{\R}{\mathbb R}
\newcommand{\C}{\mathbb C}
\newcommand{\HH}{\mathbb H}
\newcommand{\HC}{\mathbb H_{\C}}
\newcommand{\Sp}{\mathbb S}
\newcommand{\Cp}{\C_+}
\newcommand{\ii}{\mathbf i}
\newcommand{\jj}{\mathbf j}
\newcommand{\kk}{\mathbf k}
\newcommand{\D}{\mathcal D}
\newcommand{\Lap}{\Delta_4}
\newcommand{\Ind}{\mathcal I}
\newcommand{\Fu}{\mathfrak F}
\newcommand{\AM}{\mathcal{AM}}
\newcommand{\Hol}{\mathcal O}
\newcommand{\Aff}{\mathcal{Aff}_{\HH}}
\newcommand{\Per}{\operatorname{Per}}
\newcommand{\Res}{\operatorname{Res}}
\DeclareMathOperator{\Log}{Log}
\newcommand{\wind}{\operatorname{wind}}
\newcommand{\im}{\operatorname{im}}
\newcommand{\ReC}{\operatorname{Re}_{\iota}}
\newcommand{\ImC}{\operatorname{Im}_{\iota}}
\newcommand{\dd}{\,\mathrm d}
\newcommand{\T}{\mathcal T}
\newcommand{\G}{\mathcal G}
\newcommand{\Uker}{\mathcal U}
\newcommand{\Vker}{\mathcal V}
\newcommand{\pairnorm}[1]{\lVert #1\rVert^{\mathrm{ax}}}
\newcommand{\cnorm}[1]{\lvert #1\rvert_{\C}}
\newcommand{\Hone}{H_1}
\newcommand{\Z}{\mathbb Z}
\setlength{\parindent}{1.2em}
\setlength{\parskip}{2pt}
\setlength{\emergencystretch}{2em}
\setlist{itemsep=2pt,topsep=4pt}
\pagestyle{fancy}
\fancyhf{}
\fancyhead[L]{\small Affine monodromy and spherical principal parts}
\fancyhead[R]{\small\thepage}
\renewcommand{\headrulewidth}{0.3pt}
\title{\vspace{-10mm}\textbf{Affine Monodromy and Spherical Principal Parts\\in Quaternionic Fueter Inversion}\\[6pt]
\large Global obstructions, a holomorphic encoding, and sharp singularity classification}
\author{A research manuscript prepared with ChatGPT}
\date{21 September 2026}
\begin{document}
\maketitle
\thispagestyle{empty}
\begin{abstract}
Local inversion of the quaternionic Fueter map leaves a global question:
when does an axial monogenic function have a single-valued slice-regular
primitive on its original domain? We give an explicit answer by two
quaternion-valued closed one-forms. Their periods are exactly the two
coefficients of the affine monodromy of a local primitive. For a planar
base with $m$ holes, logarithmic modes split the resulting exact sequence,
and the range has real codimension $8m$. A first-order differential
expression in the monogenic field is holomorphic and equals the second
complex derivative of every local primitive. This encoding recovers all
periods and spherical principal-part coefficients by ordinary complex
contour integrals. It also exposes a constraint linking its simple- and
double-pole coefficients. At a nonreal sphere we obtain a convergent
Laurent--logarithmic normal form, an $8N$-dimensional space of singular
classes of growth $O(\rho^{-N})$, exact leading norm asymptotics, and the
sharp removability criterion $o(\rho^{-1})$, where $\rho$ is distance to
the sphere. We further give a real-axis extension theorem, explicit
same-domain counterexamples to unrestricted global surjectivity,
quantitative approximation obstructions, and period-corrected rational
approximation. The local Fueter theorem and the existence of arctangent
kernels are established background; the global obstruction and singularity
package developed here is offered as a candidate new contribution, with
its literature comparison and priority limitations stated explicitly.
\end{abstract}
\noindent\textbf{Keywords.} Quaternionic analysis; inverse Fueter map; axial
monogenic function; affine monodromy; period obstruction; spherical residue.

\medskip
\noindent\textbf{Reading conventions.} All quaternionic coefficients act
on the right. The symbol $\iota$ is a central complex imaginary unit,
not a quaternionic imaginary unit. Except in Section~\ref{sec:real}, planar
bases lie strictly above the real axis. Theorems concern axial monogenic
functions, not arbitrary solutions of the Cauchy--Fueter equation.
\newpage
\begingroup
\small
\setlength{\parskip}{0pt}
\makeatletter
\patchcmd{\l@section}{1.0em}{0.45em}{}{}
\makeatother
\tableofcontents
\endgroup
\newpage

\section{Questions, results, and relation to the supplied article}
\label{sec:intro}

\subsection{The gap to be filled}
The supplied article~\cite{Supplied} distinguishes slice regularity from
Cauchy--Fueter regularity and ends its inverse-Fueter construction with a
specific qualification: on a nonsimply connected planar set, the closed
forms used in the construction can have nonzero periods. Its inverse
therefore applies on a simply connected upper-half-plane base, with an
affine ambiguity $qa+b$. We retain precisely these conventions and develop
the missing global and singularity theory rather than repeat the classical
complex-analytic analogues in that article.

The first question is whether the two successive integrations in the
local proof conceal an additional compatibility condition. They do not:
a change of potential gives two \emph{independently defined closed forms},
both computed directly from the given monogenic field. Their periods
completely characterize global inversion. The second question is whether
every formally possible affine monodromy actually occurs. On a finitely
connected base, the answer is yes, and elementary logarithmic modes give
an explicit right inverse of the period map.

The third question concerns an isolated \emph{nonreal sphere}, rather than
an isolated point. Here the affine monodromy survives as a new part of the
principal part. We classify it together with all polar terms, recover its
coefficients from the field itself, and determine the exact growth scale.
Thus topology and singularity order become two separately measurable
features of the same function.

\subsection{Principal conclusions}
For a connected domain $U\subset\Cp=\{x+\iota r:r>0\}$, let
$g(x+Ir)=P(x,r)+IQ(x,r)$ be axial monogenic. Define
\begin{align}
 \omega_g&=\frac12(-P\dd x+Q\dd r),\label{eq:intro-omega}\\
 \theta_g&=\frac12\bigl((xP+rQ)\dd x+(rP-xQ)\dd r\bigr).
 \label{eq:intro-theta}
\end{align}
Theorem~\ref{thm:global} says that $g=\Lap f$ for a globally single-valued
slice-regular $f$ if and only if both forms have zero periods. On a loop
$\gamma$, their two periods $a_\gamma,b_\gamma$ are exactly the change
$qa_\gamma+b_\gamma$ of any local primitive.

There is also a holomorphic expression
\begin{equation}\label{eq:intro-T}
 \T g=\frac12\bigl(-P+rQ_x-\iota rP_x\bigr).
\end{equation}
It satisfies $\T(\Lap\Ind(F))=F''$ and
\begin{equation}\label{eq:intro-complexperiod}
 a_\gamma=\int_\gamma\T g\dd z,
 \qquad b_\gamma=-\int_\gamma z\T g\dd z.
\end{equation}
These complex-looking integrals actually belong to $\HH$, not just $\HC$.
Their reality is a nontrivial restriction on the possible holomorphic
encodings of monogenic fields.

For $m$ holes, the quotient of axial monogenic fields by global Fueter
images has right quaternionic dimension $2m$, or real dimension $8m$
(Theorem~\ref{thm:split}). At a sphere $u+v\Sp$, the quotient of germs of
growth $O(\rho^{-N})$ by regular germs has right quaternionic dimension
$2N$, or real dimension $8N$ (Theorem~\ref{thm:finite}). These are different
dimension counts: one is global and topological; the other is local and
filtered by singularity order.

\subsection{What is established background, and what is proposed here}
The slice-stem framework, Fueter's mapping theorem, local inverse formulas,
and the affine kernel are established mathematics
\cite{GS,GP,CSS,CPSS}. In particular, the arctangent primitives and
sphere-integrated Cauchy kernels in \cite[Theorem 5.1]{CSS} are not new.
The inverse construction in \cite[Theorem 3]{CPSS} is stated on a rectangle
in the upper half-plane. These sources are essential antecedents, not
problems purportedly solved for the first time here.

The candidate contribution is the explicit same-domain period
classification, its split range and stability statements, the holomorphic
encoding with its residue constraint, and the resulting complete
spherical principal-part and growth classification. No assertion of
absolute priority is made: the literature search described in
Section~\ref{sec:literature} did not locate this package in the inspected
sources, but it does not exclude an equivalent result elsewhere.
All mathematical assertions used for these conclusions are proved below;
their correctness does not depend on that priority assessment.

We use standard one-variable complex analysis, componentwise for
$\HC$-valued functions: Cauchy's theorem, Laurent series, analytic
continuation, and, only in Section~\ref{sec:runge}, Runge approximation.
We also use the elementary correspondence between closed-form periods
and first homology. No general sheaf-cohomology theorem or unproved
quaternionic inverse theorem is needed in the proofs.

\section{Axial fields and the Fueter map}
\label{sec:setup}

\subsection{Scalars and domains}
Write $\HH=\R+\ii\R+\jj\R+\kk\R$, with
$\ii^2=\jj^2=\kk^2=\ii\jj\kk=-1$, and put
$\Sp=\{I\in\HH:I^2=-1\}$. The complexification is
\[
 \HC=\{A+\iota B:A,B\in\HH\},\qquad \iota^2=-1,
 \quad\iota a=a\iota\quad(a\in\HH).
\]
Central conjugation is $\kappa(A+\iota B)=A-\iota B$; it is different from
quaternionic conjugation. We write $\ReC(A+\iota B)=A$,
$\ImC(A+\iota B)=B$, and
\[
 \cnorm{A+\iota B}=(|A|^2+|B|^2)^{1/2}.
\]
Multiplication by a central complex number $z$ scales this norm by $|z|$.

For a connected open $U\subset\Cp$, let
\[
 \Omega_U=\{x+Ir:x+\iota r\in U,\ I\in\Sp\}.
\]
A holomorphic $F=A+\iota B:U\to\HC$ induces
\begin{equation}\label{eq:induce}
 \Ind(F)(x+Ir)=A(x,r)+IB(x,r).
\end{equation}
Equivalently, extend $F$ to $U\cup\overline U$ by
$F(\bar z)=\kappa F(z)$ and use the usual stem construction. No condition
$F(z)\in\HH$ is imposed on the upper half-plane. In particular, the
coefficients of a Laurent series centered there belong to $\HC$.
We identify this slice-regular space with $\Hol(U;\HC)$ when convenient.

An \emph{axial field} is a function $g=P+IQ$ as in
\eqref{eq:induce}, with $P,Q:U\to\HH$ independent of $I$, not necessarily
satisfying the slice Cauchy--Riemann equations. Its pair norm is
\begin{equation}\label{eq:pairnorm}
 m_g(z)=(|P(z)|^2+|Q(z)|^2)^{1/2},\qquad
 \pairnorm{g}_K=\sup_{z\in K}m_g(z).
\end{equation}
It is equivalent to the ordinary supremum on the circularization of $K$:
\begin{equation}\label{eq:norm-equiv}
 m_g(z)\le \sup_{I\in\Sp}|P(z)+IQ(z)|\le\sqrt2\,m_g(z).
\end{equation}
The right inequality is the triangle inequality and Cauchy--Schwarz. For
the left inequality, average the squared norm over $I\in\Sp$; the cross
term has average zero.

\subsection{The differential equations}
The left Cauchy--Fueter operator and the four-dimensional Laplacian are
\[
 \D=\partial_{x_0}+\ii\partial_{x_1}+\jj\partial_{x_2}+\kk\partial_{x_3},
 \qquad \Lap=\sum_{\mu=0}^3\partial_{x_\mu}^2.
\]
Let $\AM(U)$ be the right $\HH$-vector space of smooth axial fields on
$\Omega_U$ satisfying $\D g=0$.

\begin{lemma}[Axial identities]\label{lem:axial}
For $r>0$ and $h(x+Ir)=A(x,r)+IB(x,r)$,
\begin{align}
 \D h&=A_x-B_r-2B/r+I(B_x+A_r),\label{eq:Dax}\\
 \Lap h&=A_{xx}+A_{rr}+2A_r/r\notag\\
 &\hspace{12mm}+I(B_{xx}+B_{rr}+2B_r/r-2B/r^2).
 \label{eq:Lapax}
\end{align}
Consequently $g=P+IQ$ belongs to $\AM(U)$ precisely when
\begin{equation}\label{eq:vekua}
 P_x-Q_r=2Q/r,\qquad Q_x+P_r=0.
\end{equation}
\end{lemma}
\begin{proof}
Write $v=\ii x_1+\jj x_2+\kk x_3$, $r=|v|$, and $I=v/r$.
For $e_1=\ii,e_2=\jj,e_3=\kk$,
\[
 \partial_{x_j}r=x_j/r,\qquad
 \partial_{x_j}I=e_j/r-vx_j/r^3.
\]
Thus $\sum e_j\partial_{x_j}I=-2/r$, the angular Laplacian contributes
$-2I/r^2$, and the mixed radial-angular term vanishes.
The product rule gives \eqref{eq:Dax}--\eqref{eq:Lapax}, preserving all
quaternionic factors on their displayed sides. An expression $C+IH$
vanishing for every $I$ has $C=H=0$, by comparing $I$ and $-I$; this gives
\eqref{eq:vekua}.
\end{proof}

Define the real-linear, right $\HH$-linear Fueter map on stems by
\begin{equation}\label{eq:Fu-def}
 \Fu F=\Lap\Ind(F),\qquad F\in\Hol(U;\HC).
\end{equation}
Holomorphy means $A_x=B_r$, $A_r=-B_x$. Hence
\begin{equation}\label{eq:Fu-formula}
 \Fu F=P+IQ,\qquad P=\frac{2A_r}{r},\quad
 Q=\frac{2B_r}{r}-\frac{2B}{r^2}.
\end{equation}
Equivalently,
\begin{equation}\label{eq:Fu-first}
 Q-\iota P=\frac2r F'-\frac2{r^2}\ImC F,
 \qquad
 m_{\Fu F}\le\frac2r\cnorm{F'}+\frac2{r^2}\cnorm F.
\end{equation}
This explains why, off the real axis, singularity order increases by
\emph{one} under the Fueter map, although its definition uses a
four-dimensional operator of order two.

\begin{proposition}[Fueter mapping and affine kernel]\label{prop:Fu}
The map \eqref{eq:Fu-def} takes values in $\AM(U)$, and
\begin{equation}\label{eq:kernel}
 \ker\Fu=\Aff(U):=\{z\mapsto za+b:a,b\in\HH\}.
\end{equation}
\end{proposition}
\begin{proof}
Use \eqref{eq:Fu-formula} and the Cauchy--Riemann equations. They give
\[
 P_x=2B_{rr}/r=Q_r+2Q/r,\qquad Q_x+P_r=0.
\]
Thus \eqref{eq:vekua} holds. If $P=Q=0$, put $C=B/r$.
Then $C_x=-P/2=0$ and $C_r=Q/2=0$, so $C=a\in\HH$ on connected $U$.
The Cauchy--Riemann equations imply $A_x=a,A_r=0$, whence $A=xa+b$ and
$F=za+b$. The converse follows directly from \eqref{eq:Fu-formula}.
\end{proof}

\begin{remark}
The kernel is not the set of all affine $\HC$-valued holomorphic functions.
For example, $F=\iota$ induces $I$, and $\Fu(\iota)=-2I/r^2\ne0$.
Confusing these two affine spaces would double the kernel dimension and
invalidate the period theory.
\end{remark}

\section{Canonical potentials and the global inverse theorem}
\label{sec:periods}

\begin{lemma}[Two closed forms without a prior choice of primitive]
\label{lem:forms}
For $g\in\AM(U)$, the forms $\omega_g,\theta_g$ in
\eqref{eq:intro-omega}--\eqref{eq:intro-theta} are closed. If
$\Fu F=g$ locally, $F=A+\iota B$, then
\begin{equation}\label{eq:potentials}
 C=B/r,\qquad E=A-xB/r
\end{equation}
satisfy
\begin{equation}\label{eq:potential-forms}
 \dd C=\omega_g,\qquad \dd E=\theta_g.
\end{equation}
Conversely, any such $\HH$-valued potentials produce a holomorphic stem
\begin{equation}\label{eq:reconstruct}
 F(z)=zC(z)+E(z),\qquad \Ind(F)(q)=qC(x,r)+E(x,r),
\end{equation}
with $\Fu F=g$.
\end{lemma}
\begin{proof}
For $\omega_g$, the coefficient of $\dd x\wedge\dd r$ in its exterior
derivative is $(Q_x+P_r)/2=0$. For $\theta_g$, twice this coefficient is
\[
 rP_x-Q-xQ_x-xP_r-Q-rQ_r
 =r(P_x-Q_r)-2Q-x(Q_x+P_r)=0.
\]
If $\Fu F=g$, the Cauchy--Riemann equations and
\eqref{eq:Fu-formula} give
\[
 C_x=-P/2,\quad C_r=Q/2,\quad
 E_x=(xP+rQ)/2,\quad E_r=(rP-xQ)/2.
\]
These prove \eqref{eq:potential-forms}. Conversely, set $A=xC+E$,
$B=rC$. The same four equations give
\[
 A_x=C+rQ/2=B_r,\qquad A_r=rP/2=-B_x.
\]
Thus $F$ is holomorphic, and substitution in \eqref{eq:Fu-formula}
yields $g$.
\end{proof}

The potential $E=A-xC$ is the useful change of variable. In the supplied
article the second form contains the already chosen first potential $C$.
Subtracting $\dd(xC)$ eliminates it and produces $\theta_g$ directly from
$g$. There is consequently no iterated-integral ambiguity left in the
obstruction criterion.

\begin{definition}[Affine period homomorphism]
For a piecewise $C^1$ closed curve $\gamma\subset U$, set
\begin{equation}\label{eq:Per}
 \Per_g(\gamma)=(a_\gamma,b_\gamma)
 :=\left(\int_\gamma\omega_g,\int_\gamma\theta_g\right)\in\HH^2.
\end{equation}
Integration is componentwise over the real coordinates of each
quaternion. Closedness makes this a homomorphism
$\Hone(U;\Z)\to\HH^2$.
\end{definition}

\begin{theorem}[Global inversion and affine monodromy]\label{thm:global}
For connected open $U\subset\Cp$ and $g\in\AM(U)$, the following
conditions are equivalent:
\begin{enumerate}[label=\textup{(\roman*)}]
\item There is $F\in\Hol(U;\HC)$ with $\Fu F=g$.
\item Both $\omega_g$ and $\theta_g$ are exact.
\item $\Per_g(\gamma)=(0,0)$ for every closed curve $\gamma\subset U$.
\end{enumerate}
If these conditions hold, all primitives differ by $za+b$, $a,b\in\HH$.
For a base point $z_*$, the unique primitive with $C(z_*)=E(z_*)=0$ is
\begin{equation}\label{eq:path-inverse}
 F(z)=z\int_{z_*}^{z}\omega_g+\int_{z_*}^{z}\theta_g,
\end{equation}
where either integral may be taken along any path in $U$.
Without the period-vanishing assumption, a primitive exists on the
universal cover, and its continuation around $\gamma$ changes by
\begin{equation}\label{eq:monodromy}
 F(z)\longmapsto F(z)+za_\gamma+b_\gamma.
\end{equation}
\end{theorem}
\begin{proof}
Lemma~\ref{lem:forms} proves (i)$\Rightarrow$(ii) and
(ii)$\Rightarrow$(i). Exact forms have zero periods. Conversely, when all
periods vanish, define $C,E$ by path integration from $z_*$. The difference
between any two paths is a closed curve, so both values are well defined.
On a small disk, the ordinary fundamental theorem for line integrals
shows their differentials are the prescribed forms. This proves
(iii)$\Rightarrow$(ii), and also \eqref{eq:path-inverse}.
The kernel statement is Proposition~\ref{prop:Fu}.

On the universal cover, the pulled-back forms are exact. Adding a based
loop $\gamma$ to a path adds $a_\gamma$ to $C$ and $b_\gamma$ to $E$.
Equation~\eqref{eq:reconstruct} gives \eqref{eq:monodromy}. Equivalently,
on every common continuation chart the two holomorphic branches differ
by the same affine expression. This is precisely the monodromy of the
local inverse, rather than a choice of branches of $g$, which remains
single valued.
\end{proof}

\begin{remark}[Normalization]
At a nonreal base point $z_*=x_*+\iota r_*$, the condition
$C(z_*)=E(z_*)=0$ is equivalent to $F(z_*)=0$ in $\HC$, an eight-real-
dimensional condition. It is \emph{not} equivalent to the four-real-
dimensional condition $\Ind(F)(x_*+I r_*)=0$ on one quaternionic value.
The map $(a,b)\mapsto z_*a+b$ from $\HH^2$ to $\HC$ is a real isomorphism
because $r_*>0$.
\end{remark}

The theorem gives an exact sequence, without asserting the last arrow
is surjective for arbitrary infinitely connected $U$:
\begin{equation}\label{eq:exact-general}
 0\longrightarrow\HH^2\xrightarrow{(a,b)\mapsto za+b}
 \Hol(U;\HC)\xrightarrow{\Fu}\AM(U)
 \xrightarrow{\Per}\operatorname{Hom}(\Hone(U;\Z),\HH^2).
\end{equation}
Surjectivity under explicit finite-topology hypotheses is proved in
Section~\ref{sec:split}. The homomorphism is additive even though the
coefficients are quaternions: the transition functions are \emph{added}
affine functions, not composed noncommutative fractional transformations.

\section{A holomorphic encoding of axial monogenic fields}
\label{sec:encoding}

The period forms require no derivatives of $g$. For residue extraction,
a complementary description using a single holomorphic function is more
efficient.

\begin{theorem}[The second-derivative encoding]\label{thm:T}
For every $g=P+IQ\in\AM(U)$, the function
\begin{equation}\label{eq:T}
 J_g:=\T g=\frac{-P+rQ_x-\iota rP_x}{2}
\end{equation}
is holomorphic on $U$. Every local primitive $F$ satisfies
\begin{equation}\label{eq:T-F}
 \T g=F'',\qquad
 a_\gamma=\int_\gamma J_g\dd z,\quad
 b_\gamma=-\int_\gamma zJ_g\dd z.
\end{equation}
In particular, $g$ has a global Fueter primitive if and only if both
complex integrals in \eqref{eq:T-F} vanish for every closed curve.
\end{theorem}
\begin{proof}
Let $C,E$ be local potentials. Differentiating $F=zC+E$ in $x$ gives
\begin{equation}\label{eq:Fprime}
 F'=C+\frac r2(Q-\iota P).
\end{equation}
A second $x$ derivative, with $C_x=-P/2$, is \eqref{eq:T}.
It is holomorphic because it is locally the second derivative of a
holomorphic function; the local expressions agree because
\eqref{eq:T} is defined directly from $g$.

For completeness, the Cauchy--Riemann equations of \eqref{eq:T} also
follow without a primitive. For its two components $J_0,J_1$,
\[
 (J_0)_x-(J_1)_r=\tfrac r2(Q_{xx}+P_{xr})=0,
\]
while $(J_0)_r+(J_1)_x=0$ follows from
$P_{xx}=Q_{rx}+2Q_x/r$ and $P_r=-Q_x$.

Continue a primitive along a loop. Since $\dd F'=J_g\dd z$, its first
derivative changes by $\int J_g\dd z=a_\gamma$. Also
\[
 \dd(F-zF')=-zJ_g\dd z.
\]
The affine change $za_\gamma+b_\gamma$ changes $F-zF'$ by $b_\gamma$.
This proves the second integral formula. The final assertion is
Theorem~\ref{thm:global}.
\end{proof}

\subsection{Which holomorphic functions occur?}
Let $\Hol_{\mathrm{ra}}(U;\HC)$ consist of holomorphic $J$ for which
\begin{equation}\label{eq:real-affine-condition}
 \int_\gamma J\dd z\in\HH,\qquad
 -\int_\gamma zJ\dd z\in\HH
 \quad\text{for every closed }\gamma\subset U.
\end{equation}
The subscript means that double primitives have \emph{real-quaternionic
affine monodromy}. It does not mean $J$ is real valued.

\begin{theorem}[Range and kernel of the encoding]\label{thm:T-range}
The map $\T:\AM(U)\to\Hol_{\mathrm{ra}}(U;\HC)$ is surjective. Its
kernel consists precisely of
\begin{equation}\label{eq:T-kernel}
 R_{e,h}(x+Ir)=-\frac{2e}{r}-I\frac{2(xe+h)}{r^2},
 \qquad e,h\in\HH.
\end{equation}
Each such field is itself a global Fueter image:
$R_{e,h}=\Fu\bigl(\iota(ze+h)\bigr)$.
\end{theorem}
\begin{proof}
Necessity of \eqref{eq:real-affine-condition} follows from
Theorem~\ref{thm:T}. Conversely, take a double primitive $F$ of $J$ on
the universal cover. Analytic continuation around a loop changes $F''$
by zero, hence changes $F$ by $z a+b$ with $a,b\in\HC$.
Integration of $\dd F'$ and $\dd(F-zF')$ shows that $a,b$ are the two
integrals in \eqref{eq:real-affine-condition}. They belong to $\HH$, so
Proposition~\ref{prop:Fu} implies that $\Fu F$ is invariant under every
continuation and descends to $g\in\AM(U)$. Locally $\T g=F''=J$.

If $\T g=0$, a local primitive is $F=zc+d$ with $c,d\in\HC$.
Write $c=a+\iota e,d=b+\iota h$, with all four coefficients in $\HH$.
The term $za+b$ is killed by $\Fu$, and the remaining term gives
\eqref{eq:T-kernel} by direct substitution. The parameters $e,h$ are
unique, first from the $P$ component and then from $Q$. Local parameters
therefore agree on overlaps and are global on connected $U$. The converse
and the displayed global primitive follow from the same calculation.
\end{proof}

This theorem identifies axial monogenic functions with a constrained
holomorphic function plus eight real parameters. In particular, it gives
a constructive test for proposed holomorphic encodings: arbitrary complex
residues need not satisfy \eqref{eq:real-affine-condition}.

\begin{corollary}[Singularities are detected faithfully]\label{cor:T-removable}
Let $\zeta\in\Cp$, and let $g$ be axial monogenic over a punctured disk
centered at $\zeta$. Then $g$ extends regularly across the corresponding
sphere if and only if $J_g$ extends holomorphically across $\zeta$.
\end{corollary}
\begin{proof}
One direction follows from \eqref{eq:T}. For the other, choose a
holomorphic double primitive $F_0$ of the extended $J_g$ on a disk.
The difference $g-\Fu F_0$ has zero encoding, so it is $R_{e,h}$ by
Theorem~\ref{thm:T-range}. Since the disk stays in $r>0$, this field is
smooth across its center. Both summands therefore extend.
\end{proof}

\section{Logarithmic modes realizing the affine periods}
\label{sec:modes}

Fix $\zeta=u+\iota v\in\Cp$, put $w=z-\zeta$ and $\rho=|w|$, and use
a local branch of
\begin{equation}\label{eq:logh}
 h_\zeta(z)=\frac{1}{2\pi\iota}\Log(z-\zeta).
\end{equation}
Its change around a positively oriented loop of winding one is $1$.
Define
\begin{equation}\label{eq:G-def}
 \G_{\zeta;a,b}:=\Fu\bigl(h_\zeta(z)(za+b)\bigr),\qquad a,b\in\HH.
\end{equation}
Although the stem in this expression is usually multivalued, its Fueter
image is single valued: a branch change adds $n(za+b)$, which lies in
$\ker\Fu$. Thus \eqref{eq:G-def} defines a genuine axial monogenic
function on $\Omega_{\Cp\setminus\{\zeta\}}$.

\begin{proposition}[Explicit period-normalized kernels]\label{prop:modes}
Write $\Uker_\zeta=\G_{\zeta;1,0}$ and
$\Vker_\zeta=\G_{\zeta;0,1}$. Then
$\G_{\zeta;a,b}=\Uker_\zeta a+\Vker_\zeta b$. Their scalar axial
components are
\begin{align}
 P_V&=\frac{x-u}{\pi r\rho^2},&
 Q_V&=-\frac{r-v}{\pi r\rho^2}+\frac{\log\rho}{\pi r^2},
 \label{eq:V}\\
 P_U&=\frac{x(x-u)+r(r-v)}{\pi r\rho^2}
                +\frac{\log\rho}{\pi r},&
 Q_U&=\frac{xv-ru}{\pi r\rho^2}
                +\frac{x\log\rho}{\pi r^2}.
 \label{eq:U}
\end{align}
For every loop $\gamma$ avoiding $\zeta$,
\begin{equation}\label{eq:mode-periods}
 \Per_{\G_{\zeta;a,b}}(\gamma)
       =\wind(\gamma,\zeta)(a,b).
\end{equation}
The holomorphic encoding is the rational function
\begin{equation}\label{eq:mode-T}
 \T\G_{\zeta;a,b}
 =\frac{a}{2\pi\iota(z-\zeta)}
  -\frac{\zeta a+b}{2\pi\iota(z-\zeta)^2}.
\end{equation}
\end{proposition}
\begin{proof}
Locally $h_\zeta=(\arg w-\iota\log\rho)/(2\pi)$.
Using
\[
 (\arg w)_r=(x-u)/\rho^2,\qquad
 (\log\rho)_r=(r-v)/\rho^2
\]
in \eqref{eq:Fu-formula} gives \eqref{eq:V}. For $zh_\zeta$ the
components are
\[
 A=(x\arg w+r\log\rho)/(2\pi),\qquad
 B=(r\arg w-x\log\rho)/(2\pi).
\]
Substitution gives \eqref{eq:U}; all angular terms cancel. This also
verifies directly that the formulas are branch independent.
The stem monodromy is $\wind(\gamma,\zeta)(za+b)$, so
Theorem~\ref{thm:global} proves \eqref{eq:mode-periods}.
Finally, differentiate $h_\zeta(za+b)$ twice, using
$h_\zeta'=1/(2\pi\iota w)$ and $h_\zeta''=-1/(2\pi\iota w^2)$.
Theorem~\ref{thm:T} gives \eqref{eq:mode-T}.
\end{proof}

\begin{example}[A genuine obstruction at one hole]\label{ex:annulus}
On the annulus $U=\{\varepsilon<|z-\zeta|<R\}$ with $R<v$, the field
$\Vker_\zeta$ is smooth and axial monogenic. Its period pair on a
counterclockwise circle is $(0,1)$, so it has no global slice-regular
Fueter primitive on $\Omega_U$. The field $\Uker_\zeta$ has pair $(1,0)$.
The four right multiples of either by $1,\ii,\jj,\kk$ give eight
independent real obstruction directions.
\end{example}

A vanishing first period alone is insufficient: $\Vker_\zeta$ supplies
an explicit counterexample. Thus a criterion based solely on the first
integration in the supplied article would miss half the obstruction.

\section{Finite topology: an exact sequence with an explicit splitting}
\label{sec:split}

For the finite-rank statements, assume $\Hone(U;\Z)$ is freely generated
by $m$ loops $\gamma_1,\ldots,\gamma_m$, and choose points
$\zeta_1,\ldots,\zeta_m\in\Cp\setminus U$ with
\begin{equation}\label{eq:winding-basis}
 \wind(\gamma_j,\zeta_k)=\delta_{jk}.
\end{equation}
This includes bounded finitely connected Jordan domains whose closure
lies in $\Cp$, as well as disks with finitely many punctures. In the
Jordan case, choose one point in each bounded complementary component
and a positively oriented loop surrounding that hole. Loops are oriented
counterclockwise around holes, not with the boundary orientation of $U$.

Write $\Per g=((a_j,b_j))_{j=1}^m$. Define a linear synthesis map
\begin{equation}\label{eq:S}
 S((a_j,b_j))=\sum_{j=1}^m\G_{\zeta_j;a_j,b_j}.
\end{equation}

\begin{theorem}[Split global range theorem]\label{thm:split}
Under \eqref{eq:winding-basis}, the sequence
\begin{equation}\label{eq:exact-finite}
 0\longrightarrow\HH^2\longrightarrow\Hol(U;\HC)
 \xrightarrow{\Fu}\AM(U)\xrightarrow{\Per}\HH^{2m}
 \longrightarrow0
\end{equation}
is exact. The map $S$ satisfies $\Per\circ S=\mathrm{id}$, and
\begin{equation}\label{eq:projection}
 \Pi=\mathrm{id}_{\AM(U)}-S\Per
\end{equation}
is a projection onto $\im\Fu$. In particular,
\begin{equation}\label{eq:decomp}
 \AM(U)=\im\Fu\ \oplus\
 \operatorname{span}_{\HH}\{\Uker_{\zeta_j},\Vker_{\zeta_j}:1\le j\le m\},
\end{equation}
and
\begin{equation}\label{eq:codim}
 \dim_{\HH}(\AM(U)/\im\Fu)=2m,\qquad
 \dim_{\R}(\AM(U)/\im\Fu)=8m.
\end{equation}
Every $g$ has a decomposition
\[
 g=\Fu F+\sum_j(\Uker_{\zeta_j}a_j+\Vker_{\zeta_j}b_j),
\]
unique after the normalization $F(z_*)=0$ in $\HC$.
\end{theorem}
\begin{proof}
Theorem~\ref{thm:global} proves exactness up to the last term. Equations
\eqref{eq:mode-periods} and \eqref{eq:winding-basis} give
$\Per S=\mathrm{id}$, proving surjectivity and the splitting. Therefore
$\Per\Pi=0$ and $\Pi$ is the identity on $\ker\Per=\im\Fu$; hence
$\Pi^2=\Pi$. A linear combination of the displayed modes belonging to
$\im\Fu$ has all periods zero, so each coefficient vanishes. This proves
the direct sum and the dimension counts. Finally, the residual field
$\Pi g$ has a primitive by Theorem~\ref{thm:global}, and the specified
normalization removes its affine ambiguity.
\end{proof}

The complement in \eqref{eq:decomp} depends on the chosen hole points and
homology basis. The quotient and its dimension do not. Moving a hole
point inside the same complementary component changes the corresponding
mode by a zero-period field, hence by a global Fueter image. This is the
precise sense in which the obstruction space is canonical while its
explicit representatives are not.

\section{Stability and the cost of ignoring periods}
\label{sec:stability}

\begin{proposition}[Period bounds and a distance obstruction]
\label{prop:bounds}
Let $\gamma$ have length $L_\gamma$ and put
$R_\gamma=\sup_{z\in\gamma}|z|$. Then
\begin{equation}\label{eq:period-bounds}
 |a_\gamma|\le\frac{L_\gamma}{2}\pairnorm g_\gamma,
 \qquad
 |b_\gamma|\le\frac{L_\gamma R_\gamma}{2}\pairnorm g_\gamma.
\end{equation}
If a compact $K\subset U$ contains the finite basis loops and
$h\in\im\Fu$, then
\begin{equation}\label{eq:distance}
 \pairnorm{g-h}_K\ge
 \max_{1\le j\le m}
 \left\{\frac{2|a_j|}{L_j},\frac{2|b_j|}{L_jR_j}\right\}.
\end{equation}
\end{proposition}
\begin{proof}
For a parametrized curve $(x(t),r(t))$, the integrand of $\omega_g$ has
norm at most $m_g\sqrt{x'^2+r'^2}/2$. For $\theta_g$, its coefficients
of $P,Q$ are $xx'+rr'$ and $rx'-xr'$, whose squared sum is
$(x^2+r^2)(x'^2+r'^2)$. Quaternion-valued Cauchy--Schwarz therefore gives
\eqref{eq:period-bounds} after integration. Apply these estimates to
$g-h$, whose periods equal those of $g$, to obtain \eqref{eq:distance}.
\end{proof}

Thus a nonzero-period field cannot even be uniformly approximated by
global Fueter images on a compact set carrying its obstructing cycles.
This is not merely a failure of a particular integral formula. It is a
positive lower bound on the error of \emph{every} such approximation.

\begin{corollary}[Closed range and bounded finite-rank correction]
\label{cor:closed}
The range $\im\Fu$ is closed in $\AM(U)$ for locally uniform convergence
in the pair norm, on every connected $U\subset\Cp$. Under the hypotheses
of Theorem~\ref{thm:split}, $\Pi$ is continuous and
\begin{equation}\label{eq:correction-bound}
 \pairnorm{g-\Pi g}_K\le\sum_j
 \left(\pairnorm{\Uker_{\zeta_j}}_K|a_j|
       +\pairnorm{\Vker_{\zeta_j}}_K|b_j|\right).
\end{equation}
\end{corollary}
\begin{proof}
Each period is a continuous linear functional on the compact-open space,
by \eqref{eq:period-bounds}. Theorem~\ref{thm:global} identifies the range
as the intersection of their kernels, which is closed. In the finite
case, $S$ is continuous because its image is finite dimensional, and the
triangle inequality gives \eqref{eq:correction-bound}.
\end{proof}

\begin{proposition}[A quantitative normalized inverse]\label{prop:inverse-est}
Suppose $g$ has zero periods, and $F$ is normalized by $F(z_*)=0$.
Let $K\subset U$ be compact. Choose a compact $L\subset U$ containing
paths from $z_*$ to every point of $K$, all of length at most $\ell$,
and put $R_L=\sup_L|z|$. Then
\begin{align}
 \cnorm{F(z)}&\le\frac{\ell}{2}(|z|+R_L)\pairnorm g_L,
                 &&z\in K,\label{eq:inverse-est}\\
 \cnorm{F'(z)}&\le\frac{\ell+r}{2}\pairnorm g_L,
                 &&z=x+\iota r\in K.\label{eq:inverse-der-est}
\end{align}
Such $L$ and $\ell$ exist for every compact $K$.
\end{proposition}
\begin{proof}
Path integration gives $|C|\le\ell\pairnorm g_L/2$ and
$|E|\le\ell R_L\pairnorm g_L/2$. Use $F=zC+E$ and
\eqref{eq:Fprime}. To obtain uniformly bounded paths, cover $K$ by
finitely many disks relatively compact in $U$, join their centers to
$z_*$ by finitely many polygonal paths, and add radial segments within
the disks. The union can be placed in a compact subset of $U$.
\end{proof}

These estimates supply an explicit continuous inverse after affine
normalization, rather than just an abstract open-mapping argument.
They also specify which data are needed: a compact carrying the periods
for the correction, and a compact carrying reconstruction paths for the
primitive.

\section{Crossing the real axis and genuine slice domains}
\label{sec:real}

The obstructions are not artifacts of product domains that avoid the
real axis. Let $D\subset\C$ be connected, open, conjugation invariant,
and meeting $\R$, and let $\Omega_D$ be its circularization. By an axial
monogenic field on $\Omega_D$ we mean a smooth Cauchy--Fueter solution
whose signed axial components satisfy $P(x,-r)=P(x,r)$ and
$Q(x,-r)=-Q(x,r)$. Such components are obtained from opposite points of
any fixed slice; their parity makes the description consistent at the
real axis.

\begin{lemma}[Connected upper part and local inversion at a real point]
\label{lem:real-local}
The upper part $U=D\cap\Cp$ is connected. Every axial monogenic field
near a real point has a local slice-regular Fueter primitive there.
\end{lemma}
\begin{proof}
A path in $D$ between upper-half-plane points can be folded by
$x+\iota r\mapsto x+\iota|r|$. Conjugation invariance keeps the folded
path in $D$. Its compact image has positive distance from $\C\setminus D$;
a sufficiently small upward translate lies entirely in $D\cap\Cp$.
Short vertical segments adjust its endpoints. Thus $U$ is path connected.

For the second assertion, take a small symmetric disk about the real
point. The signed functions $P,Q$ are smooth. The forms $\omega_g$ and
$\theta_g$ have no denominators and extend smoothly across $r=0$.
They are closed away from that line by Lemma~\ref{lem:forms}, and hence
closed on the whole disk by continuity. Both forms are invariant under
the pullback of $r\mapsto-r$. Their primitives $C,E$, normalized at the
real center, are therefore even in $r$. Set $A=xC+E$, $B=rC$.
The proof of Lemma~\ref{lem:forms} proves the Cauchy--Riemann equations
for $r\ne0$, and continuity proves them on the real line. Thus
$A+\iota B$ is a holomorphic stem with the required symmetry.
Its Fueter image equals $g$ away from the axis and hence everywhere.
The induced function is smooth at real points, for its holomorphic Taylor
coefficients at a real center are in $\HH$.
\end{proof}

\begin{theorem}[Same-domain inversion on slice domains]\label{thm:real}
An axial monogenic $g$ on $\Omega_D$ has a global slice-regular primitive
on \emph{that same domain} if and only if its two period forms vanish
on all loops in $U=D\cap\Cp$. The primitive is unique up to $qa+b$,
$a,b\in\HH$.
\end{theorem}
\begin{proof}
Necessity follows by restriction and Theorem~\ref{thm:global}.
For sufficiency, construct a global upper stem $F$ on $U$. At a real
point choose the local symmetric primitive $F_0$ of
Lemma~\ref{lem:real-local}. On the upper half of its disk,
$F-F_0=za+b$ with $a,b\in\HH$, by Proposition~\ref{prop:Fu}.
Thus $F$ extends across that real interval as the symmetric holomorphic
stem $F_0+za+b$. The local extensions agree on overlaps by holomorphic
uniqueness. Reflection defines the lower half, and these extensions glue
to a single stem on $D$. The kernel statement follows from the connected
upper part.
\end{proof}

\subsection{An explicit obstruction on \texorpdfstring{$\HH$}{H} minus a sphere}
Fix $\zeta=u+\iota v$, $v>0$, and locally put
\begin{equation}\label{eq:reflected-log}
 H_\zeta(z)=\frac1{2\pi\iota}
      \bigl(\Log(z-\zeta)-\Log(z-\bar\zeta)\bigr),
 \qquad H_\zeta'(z)=\frac{v}{\pi((z-u)^2+v^2)}.
\end{equation}
Near every real point a branch can be chosen real on the real line,
since the logarithm's argument has modulus one there. Thus it is locally
an intrinsic stem. Its continuation about $\zeta$ adds $1$, and about
$\bar\zeta$ adds $-1$. Consequently
\begin{equation}\label{eq:reflected-modes}
 \widehat\G_{\zeta;a,b}
   :=\Fu\bigl(H_\zeta(z)(za+b)\bigr)
\end{equation}
is a global smooth axial monogenic field on
$\HH\setminus(u+v\Sp)$, including the real axis.

\begin{corollary}[Unrestricted global surjectivity fails]
\label{cor:counterexample}
For $(a,b)\ne(0,0)$, the field \eqref{eq:reflected-modes} has no
single-valued slice-regular Fueter primitive on
$\HH\setminus(u+v\Sp)$. Its period pair on a small upper loop about
$\zeta$ is $(a,b)$.
\end{corollary}
\begin{proof}
On the upper half-plane, $H_\zeta-h_\zeta$ is single-valued holomorphic:
$\Log(z-\bar\zeta)$ has a global branch there. Therefore the reflected
mode differs from $\G_{\zeta;a,b}$ by a global Fueter image on the upper
base. It has the same periods, and Theorem~\ref{thm:real} applies.
\end{proof}

For $u=0,v=1$, $H_\zeta$ is locally $\pi^{-1}\arctan z$ plus a real
constant. The associated arctangent kernels are already present in
\cite[Theorem 5.1]{CSS}; the point here is their explicit obstruction to
a \emph{single-valued primitive on the entire original punctured domain}.
There is no conflict with local inversion, with a primitive on a cut
domain, or with the rectangle theorem of \cite{CPSS}. Any global use of
an inverse formula involving an arctangent or logarithm must account for
these distinctions.

\begin{corollary}[Which holes matter when the real axis is present]
\label{cor:real-holes}
Consider a conjugation-invariant disk from which finitely many disjoint
closed disks compactly contained in it have been removed, with each removed disk either centered
on the real axis or paired with its conjugate away from that axis.
If there are $m$ nonreal conjugate pairs, the same-domain Fueter range
among smooth axial monogenic fields has real codimension $8m$.
\end{corollary}
\begin{proof}
The upper part has $m$ holes; the removed real-centered disks meet its
boundary and contribute no upper homology generator. Theorem~\ref{thm:real}
gives the kernel of its period map. The reflected modes
\eqref{eq:reflected-modes}, one for each upper hole, are smooth throughout
the original circular domain and realize arbitrary period pairs.
The proof of Theorem~\ref{thm:split} now applies verbatim.
\end{proof}

\section{Spherical residues and a Laurent--logarithmic normal form}
\label{sec:singularities}

Fix $\zeta=u+\iota v\in\Cp$ and $0<\varepsilon<v$. Write
\[
 U^*=\{0<|z-\zeta|<\varepsilon\},\qquad S_\zeta=u+v\Sp.
\]
For $q=x+Ir$ near this sphere,
\begin{equation}\label{eq:distance-sphere}
 \operatorname{dist}(q,S_\zeta)
 =\sqrt{(x-u)^2+(r-v)^2}=|z-\zeta|=\rho.
\end{equation}
Indeed, the closest point on the sphere has the same imaginary direction
as $q$, because that maximizes the relevant Euclidean inner product.
A function is \emph{regular across the sphere} if it extends to an axial
monogenic field on the circularization of the full disk.

\begin{definition}[Affine spherical residues]\label{def:residues}
For $g\in\AM(U^*)$, define
\begin{equation}\label{eq:res-pair}
 a(g;\zeta)=\int_{|z-\zeta|=s}\omega_g,\qquad
 b(g;\zeta)=\int_{|z-\zeta|=s}\theta_g,
 \quad 0<s<\varepsilon,
\end{equation}
with positive orientation. Closedness makes these independent of $s$.
They are the two coefficients of the affine monodromy, not the ordinary
three-dimensional flux residue of an isolated point singularity.
\end{definition}

\begin{theorem}[Full spherical normal form]\label{thm:laurent-log}
Every $g\in\AM(U^*)$ has a unique representation
\begin{equation}\label{eq:full-normal}
 g=g_{\mathrm{reg}}+\G_{\zeta;a,b}
       +\sum_{k=1}^{\infty}\Fu\bigl((z-\zeta)^{-k}c_k\bigr),
\end{equation}
where $a,b\in\HH$ are \eqref{eq:res-pair}, $c_k\in\HC$,
$g_{\mathrm{reg}}$ extends regularly across the sphere, and the series
and all of its derivatives converge uniformly on compact subannuli.
The negative stem series $\sum_{k\ge1}(z-\zeta)^{-k}c_k$ converges
there as a holomorphic $\HC$-valued function.
\end{theorem}
\begin{proof}
Subtract $\G_{\zeta;a,b}$. The residual $g_0$ has zero periods on the
punctured disk, since its first homology is generated by one circle.
Theorem~\ref{thm:global} supplies a single-valued holomorphic primitive
$F_0$. Its componentwise Laurent expansion splits as
\[
 F_0(z)=F_+(z)+\sum_{k\ge1}(z-\zeta)^{-k}c_k,
\]
where $F_+$ is holomorphic on the full disk. Apply $\Fu$ and set
$g_{\mathrm{reg}}=\Fu F_+$. Laurent series and all their complex
derivatives converge uniformly on compact subannuli. Formula
\eqref{eq:Fu-formula}, and smoothness of the axial coordinates for
$r>0$, give the asserted convergence after applying $\Fu$ and taking
any further real derivatives.

For uniqueness, periods first determine $a,b$. Suppose a negative stem
series $H$ has Fueter image regular across the sphere. That extended
field has a holomorphic primitive $K$ on the full disk by
Theorem~\ref{thm:global}. On the punctured disk,
$\Fu(H-K)=0$, so $H-K=za+b_0$ with $a,b_0\in\HH$.
The right side and $K$ are holomorphic through $\zeta$; hence every
negative Laurent coefficient of $H$ is zero. This proves uniqueness of
all $c_k$, and then of $g_{\mathrm{reg}}$.
\end{proof}

The logarithmic term is not an optional rewriting of the polar part.
Every displayed polar Fueter term has a single-valued stem and hence
zero affine residues. No sum of such terms can represent a nonzero
$\G_{\zeta;a,b}$ on an annulus carrying the period cycle.

\subsection{Residues and all multipoles directly from the field}
\begin{theorem}[Contour extraction and the residue constraint]
\label{thm:res-encoding}
Let $J_g=\T g$ have Laurent expansion
$J_g(z)=\sum_{n\in\Z}j_n(z-\zeta)^n$. Then
\begin{align}
 a&=2\pi\iota j_{-1},&
 b&=-2\pi\iota\bigl(j_{-2}+\zeta j_{-1}\bigr),
 \label{eq:res-J}\\
 c_k&=\frac{j_{-k-2}}{k(k+1)}
 =\frac1{2\pi\iota k(k+1)}
    \int_{|z-\zeta|=s}(z-\zeta)^{k+1}J_g(z)\dd z,
 &&k\ge1.\label{eq:ck-J}
\end{align}
The two lowest negative coefficients obey the exact condition
\begin{equation}\label{eq:res-constraint}
 \boxed{\displaystyle
 j_{-1}=\frac{\iota}{v}\ReC(j_{-2}).}
\end{equation}
Conversely, a holomorphic $J$ on $U^*$ is $J_g$ for some axial
monogenic $g$ if and only if \eqref{eq:res-constraint} holds.
In particular, $J_g$ cannot have an isolated pole of order exactly one.
\end{theorem}
\begin{proof}
The integral formulas in Theorem~\ref{thm:T} and ordinary Laurent
coefficient extraction give \eqref{eq:res-J}. Since $a,b\in\HH$, write
$j_{-1}=\iota\alpha$ with $\alpha\in\HH$, as forced by the first
formula. Write $j_{-2}=c+\iota d$, $c,d\in\HH$.
Then
\[
 j_{-2}+\zeta j_{-1}=(c-v\alpha)+\iota(d+u\alpha).
\]
The second expression in \eqref{eq:res-J} belongs to $\HH$ precisely
when $c=v\alpha$, proving \eqref{eq:res-constraint}.

For the polar coefficients, apply $\T$ to \eqref{eq:full-normal}.
The regular part has holomorphic encoding, the logarithmic mode has
encoding \eqref{eq:mode-T}, and
\[
 \T\Fu\bigl((z-\zeta)^{-k}c_k\bigr)
 =k(k+1)(z-\zeta)^{-k-2}c_k.
\]
This proves \eqref{eq:ck-J}. Conversely, on a punctured disk every loop
is homologous to a multiple of one circle. Condition
\eqref{eq:res-constraint} is exactly the requirement that both integrals
in \eqref{eq:real-affine-condition} belong to $\HH$.
Theorem~\ref{thm:T-range} produces the required field.
Finally, if $j_{-2}=0$ then \eqref{eq:res-constraint} gives $j_{-1}=0$,
ruling out a simple pole.
\end{proof}

The double-pole coefficient is free in $\HC$, but its real-quaternionic
part forces the simple-pole coefficient. More explicitly, if
$j_{-2}=c+\iota d$, then
\begin{equation}\label{eq:res-from-double}
 a=-\frac{2\pi}{v}c,\qquad
 b=2\pi d+\frac{2\pi u}{v}c.
\end{equation}
Thus the two quaternionic residues are equivalently one unrestricted
complexified quaternion. Formula \eqref{eq:ck-J} recovers every higher
multipole without choosing, computing, or normalizing a Fueter primitive.

\section{Sharp growth classification at a nonreal sphere}
\label{sec:growth}

\begin{lemma}[Growth of a single-valued primitive]\label{lem:growth-inverse}
Suppose $g_0\in\AM(U^*)$ has zero periods and
$m_{g_0}(z)=O(\rho^{-N})$ for an integer $N\ge1$.
Every single-valued holomorphic primitive $F_0$ satisfies
\begin{equation}\label{eq:primitive-growth}
 \cnorm{F_0(z)}=
 \begin{cases}
 O(1+|\log\rho|),&N=1,\\
 O(\rho^{1-N}),&N\ge2.
 \end{cases}
\end{equation}
\end{lemma}
\begin{proof}
Use the global potentials $C,E$ of Lemma~\ref{lem:forms}.
On a sufficiently small disk, $|z|$ is bounded, so
\[
 |\nabla C|\le c_1\rho^{-N},\qquad
 |\nabla E|\le c_2\rho^{-N}.
\]
Fix an outer circle inside $U^*$. To reach a point $\zeta+\rho e^{\iota t}$,
first follow that circle from a fixed base point to angle $t$, and then
follow the radial segment. The first path contributes a uniformly bounded
amount; the second contributes at most a constant times
$\int_\rho^{s_0}s^{-N}\dd s$.
This is $O(1+|\log\rho|)$ for $N=1$ and $O(\rho^{1-N})$ for $N\ge2$.
The estimates are uniform in $t$. Since $F_0=zC+E$ and $z$ remains
bounded, \eqref{eq:primitive-growth} follows. Adding an affine kernel
term does not change these bounds.
\end{proof}

\begin{theorem}[Finite-order principal parts and their dimension]
\label{thm:finite}
For $N\ge1$, a field $g\in\AM(U^*)$ satisfies
$m_g(z)=O(\rho^{-N})$ if and only if its unique normal form is
\begin{equation}\label{eq:finite-normal}
 g=g_{\mathrm{reg}}+\G_{\zeta;a,b}
       +\sum_{k=1}^{N-1}\Fu\bigl((z-\zeta)^{-k}c_k\bigr).
\end{equation}
Here $a,b\in\HH$ and all $c_k\in\HC$ are arbitrary and independent.
If $\mathscr M_N$ is the space of germs satisfying this growth bound,
and $\mathscr M_0$ the space of regular germs, then
\begin{equation}\label{eq:local-dim}
 \mathscr M_N/\mathscr M_0\ \cong\
 \HH^2\oplus\HC^{\,N-1},\qquad
 \dim_\R(\mathscr M_N/\mathscr M_0)=8N.
\end{equation}
The isomorphism is one of right $\HH$-vector spaces; its quaternionic
dimension is $2N$.
\end{theorem}
\begin{proof}
The logarithmic modes \eqref{eq:V}--\eqref{eq:U} have growth
$O(\rho^{-1})$ because $r$ is bounded away from zero. Subtracting the
residue mode preserves the assumed $O(\rho^{-N})$ bound and gives
zero periods. Choose a single-valued primitive $F_0$.
Lemma~\ref{lem:growth-inverse} and Laurent's coefficient formula show
that its negative coefficients of order $k\ge N$ vanish. Indeed their
norm is at most
\[
 \rho^k\sup_{|z-\zeta|=\rho}\cnorm{F_0(z)},
\]
which tends to zero for every $k\ge1$ when $N=1$, and for every
$k\ge N$ when $N\ge2$. This proves \eqref{eq:finite-normal}.

Conversely, \eqref{eq:Fu-first} gives
$\Fu((z-\zeta)^{-k}c_k)=O(\rho^{-k-1})$, while the regular part is
bounded and the logarithmic mode is $O(\rho^{-1})$. Thus every
\eqref{eq:finite-normal} has the stated bound.
Uniqueness is Theorem~\ref{thm:laurent-log}, and every choice of the
coefficients constructs a valid field, proving the isomorphism.
Finally $\dim_\R\HH^2=8$ and $\dim_\R\HC=8$.
\end{proof}

\subsection{Exact leading size, not just an upper bound}
\begin{theorem}[Leading norm asymptotics]\label{thm:asymptotic}
For a nonremovable finite-order singularity, let $N$ be its least
integer growth order. If $N=1$, then $(a,b)\ne(0,0)$ and
\begin{equation}\label{eq:critical-asymptotic}
 m_g(z)=\frac{\cnorm{\zeta a+b}}{\pi v\rho}
             +O(1+|\log\rho|).
\end{equation}
If $N\ge2$, the top coefficient $c_{N-1}$ is nonzero and
\begin{equation}\label{eq:pole-asymptotic}
 m_g(z)=\frac{2(N-1)\cnorm{c_{N-1}}}{v\rho^N}
             +O(\rho^{1-N}).
\end{equation}
The estimates are uniform in the polar angle of $z-\zeta$.
Equivalently, $J_g$ has a pole of order exactly $N+1$, and
\begin{equation}\label{eq:J-asymptotic}
 m_g(z)\sim \frac{2\cnorm{j_{-N-1}}}{vN}\rho^{-N}.
\end{equation}
\end{theorem}
\begin{proof}
For $N\ge2$, let $k=N-1$ and $F=w^{-k}c_k$.
Equation~\eqref{eq:Fu-first} gives
\[
 Q-\iota P=-\frac{2k}{v}w^{-k-1}c_k+O(\rho^{-k}).
\]
Here replacing $r$ by $v$ costs $O(\rho^{-k})$, and the term involving
$\ImC F/r^2$ has that same bound. Central complex multiplication scales
$\cnorm{\cdot}$, while $\cnorm{Q-\iota P}=m_g$. All lower polar terms,
the residue modes, and the regular term contribute at most
$O(\rho^{-k})$. This proves \eqref{eq:pole-asymptotic}.

For $N=1$, use a branch of $F=h_\zeta(za+b)$ whose argument is bounded
at the point being estimated. Uniformly in angle,
\[
 F'=\frac{\zeta a+b}{2\pi\iota w}+O(1+|\log\rho|),
 \qquad \cnorm F=O(1+|\log\rho|).
\]
The branch may vary with the point, since its change is a real-quaternionic
affine term killed by $\Fu$. Substitution in \eqref{eq:Fu-first} gives
\eqref{eq:critical-asymptotic}. Its leading constant is positive:
\[
 \cnorm{\zeta a+b}^2=|ua+b|^2+v^2|a|^2,
\]
which vanishes only if $a=b=0$.
The formulas for the encoding of a mode and a polar term show that its
top coefficient is $-(\zeta a+b)/(2\pi\iota)$ when $N=1$, and
$N(N-1)c_{N-1}$ when $N\ge2$. They prove
\eqref{eq:J-asymptotic} and the assertion about pole order.
\end{proof}

\begin{corollary}[Sharp removability and the critical residue criterion]
\label{cor:sharp}
If $m_g(z)=o(\rho^{-1})$ uniformly as $\rho\downarrow0$, then $g$
extends regularly across $S_\zeta$. The assertion is false with
$O(\rho^{-1})$ in place of $o(\rho^{-1})$.
Under the bound $O(\rho^{-1})$, removability is equivalent to
$a(g;\zeta)=b(g;\zeta)=0$.
\end{corollary}
\begin{proof}
Under the critical big-$O$ bound, Theorem~\ref{thm:finite} gives
$g=g_{\mathrm{reg}}+\G_{\zeta;a,b}$. Thus zero residues imply
removability, and nonzero residues obstruct it.
Under the little-$o$ bound, the circle estimates
\eqref{eq:period-bounds}, with length $2\pi\rho$ and bounded $|z|$,
force both constant periods to be zero. Alternatively,
\eqref{eq:critical-asymptotic} excludes a nonzero residue mode.
The field $\Vker_\zeta$ proves sharpness.
\end{proof}

\begin{corollary}[No fractional finite singularity orders]
\label{cor:fractional}
For a real exponent $\alpha\ge0$, a bound $m_g=O(\rho^{-\alpha})$
allows precisely the principal parts in \eqref{eq:finite-normal} with
$N=\lfloor\alpha\rfloor$, interpreting $N=0$ as no principal part.
The corresponding singular quotient has real dimension
$8\lfloor\alpha\rfloor$. In particular, every nonremovable singularity
of finite power growth has a positive integer growth order.
\end{corollary}
\begin{proof}
Apply Theorem~\ref{thm:finite} with the integer
$\max(1,\lceil\alpha\rceil)$. If a surviving principal part had least
integer order $N>\alpha$, Theorem~\ref{thm:asymptotic} would contradict
the assumed bound. The converse follows from the same theorem and the
boundedness of regular germs.
\end{proof}

\subsection{What zero residues do and do not say}
At order $N\ge2$, zero affine residues no longer imply removability.
For instance,
\[
 g=\Fu((z-\zeta)^{-1})
\]
has zero periods, holomorphic encoding $2(z-\zeta)^{-3}$, and
$m_g\sim2/(v\rho^2)$. At infinite order,
\[
 g=\Fu\bigl(\exp((z-\zeta)^{-1})\bigr)
\]
also has zero periods but has an essential holomorphic encoding.
It therefore admits no finite power bound. These examples separate
three notions: a global single-valued primitive, vanishing affine
residues, and removability. The first two are equivalent on a punctured
disk; neither implies the third without a growth hypothesis.

\section{Period-corrected rational approximation}
\label{sec:runge}

Assume the finite-topology setting of Section~\ref{sec:split}, with one
chosen pole point in each bounded complementary component and with
$\infty$ also allowed as a pole. For this section assume, as in the
usual finite-connectivity Runge setting, that these are all components
of the complement in the Riemann sphere. Let $\mathcal R$ be the
$\HC$-valued rational functions with poles only at the chosen points
and at infinity. These are \emph{upper-stem rational functions}; after
reflection to the lower component their induced slice functions need
not be rational functions with coefficients in $\HH$ on a connected
symmetric planar domain.

\begin{theorem}[Approximation with exactly preserved periods]
\label{thm:runge}
For every $g\in\AM(U)$ there exist $R_n\in\mathcal R$ such that
\begin{equation}\label{eq:runge}
 g_n=\Fu R_n+S(\Per g)\longrightarrow g
\end{equation}
locally uniformly, together with all real derivatives on compact subsets
of $\Omega_U$. Every approximant has exactly the same affine periods as
$g$. The $2m$ quaternionic modes in $S$ are minimal: a fixed finite-
dimensional real correction space whose addition to global Fueter
images is dense in $\AM(U)$ must have dimension at least $8m$.
\end{theorem}
\begin{proof}
By Theorem~\ref{thm:split}, write $g=\Fu F+S(\Per g)$ with $F$
holomorphic. Classical Runge approximation, applied to its four complex
coordinates, gives rational $R_n$ with the prescribed pole set and
$R_n\to F$ uniformly on compact subsets of $U$. Cauchy estimates on
slightly larger compact sets give convergence of all complex
derivatives. Formula~\eqref{eq:Fu-formula} and the smooth axial
coordinate map away from $r=0$ give \eqref{eq:runge} with all real
derivatives. Each $\Fu R_n$ has zero periods, proving the exact
preservation assertion.

Let $W$ be a finite-dimensional real correction space with the stated
density property. The continuous surjective period map takes
$\im\Fu+W$ into $\Per(W)$. This image is a finite-dimensional subspace
of $\HH^{2m}$, hence closed. Density forces
$\Per(W)=\HH^{2m}$, so $\dim_\R W\ge8m$.
\end{proof}

This theorem transfers a standard scalar approximation theorem through
an explicitly corrected Fueter inverse. The Runge input itself is not a
novel claim. The additional information is the obstruction, its exact
preservation at every stage, and the sharp number of correction modes.

\section{Worked calculations and a reconstruction procedure}
\label{sec:examples}

\subsection{One annulus, two independent obstructions}
Take $\zeta=\iota$ and
$U=\{1/4<|z-\iota|<1/2\}$. For the loop $|z-\iota|=s$,
\[
 \Per\Uker_{\iota}=(1,0),\qquad
 \Per\Vker_{\iota}=(0,1).
\]
The field $g=\Uker_{\iota}\ii+\Vker_{\iota}\jj$ has monodromy
$q\ii+\jj$. Its holomorphic encoding is
\[
 J_g(z)=\frac{\ii}{2\pi\iota(z-\iota)}
       -\frac{\iota\ii+\jj}{2\pi\iota(z-\iota)^2}.
\]
In this formula $\iota\ii$ is a central complex multiple of the
quaternionic unit $\ii$; it is not $-1$.
For $s=1/3$, every global Fueter image $h$ satisfies
\[
 \pairnorm{g-h}_{|z-\iota|=1/3}\ge\frac3\pi,
\]
by the first-period bound. Thus the obstruction remains visible to any
approximation that samples an entire loop, irrespective of how many
terms are used in the candidate primitive.

\subsection{A second-order singularity with prescribed residues}
Set $c=1+\iota\kk\in\HC$ and
\[
 g=\G_{\iota;\ii,\jj}+\Fu((z-\iota)^{-1}c).
\]
Then
\[
 J_g=\frac{\ii}{2\pi\iota(z-\iota)}
     -\frac{\iota\ii+\jj}{2\pi\iota(z-\iota)^2}
     +\frac{2(1+\iota\kk)}{(z-\iota)^3}.
\]
The simple and double coefficients recover $(a,b)=(\ii,\jj)$; the
third-order coefficient divided by $2$ recovers $c$. Since
$\cnorm c=\sqrt2$ and $v=1$, the field has
$m_g\sim2\sqrt2\,\rho^{-2}$. Subtracting its residues alone leaves a
nonremovable second-order singularity. Subtracting the extracted polar
term as well leaves a field regular across the sphere.

\subsection{An implementation-oriented procedure}
For a field given by evaluable axial components $P,Q$ on a finitely
connected $U$, the theory gives the following concrete procedure.
First verify the two Vekua equations, or determine that they are assumed
as exact input. Fix oriented homology loops and one point in each hole.
Compute their two integrals in \eqref{eq:Per}. Form the explicit
correction $S(\Per g)$ using \eqref{eq:V}--\eqref{eq:U}. The residual has
zero periods. Finally, integrate its two closed forms from a base point
and use \eqref{eq:path-inverse}. This produces a normalized global
primitive of the residual and a complete certificate for whether the
uncorrected field has a global primitive.

When an isolated spherical singularity is present, compute $J_g$ from
\eqref{eq:T} and use \eqref{eq:res-J}--\eqref{eq:ck-J} for its
coefficients. These contour formulas are useful even when symbolic
integration of the two potentials is difficult. Derivatives of noisy
data can be unstable, so the derivative-free period forms are preferable
for measuring just the affine obstruction. The error bounds in
Section~\ref{sec:stability} quantify perturbations of those periods.

The accompanying \texttt{verify.py} checks the kernel equations and
encoding identities symbolically and checks representative loop
integrals numerically. The proofs in this paper are independent of that
script. Numerical near-zero periods alone do not certify exact
vanishing for arbitrary input data.

\section{Literature comparison, verification status, and further questions}
\label{sec:literature}

\subsection{Comparison with the sources examined}
The attached exposition~\cite{Supplied} supplies the local two-integration
construction and identifies possible periods as a reason not to claim
global inversion. Our starting substitution $E=A-xB/r$ removes the
nested potential and makes those periods independently computable.
The derivative identities and affine kernel are retained, not
reclassified as original results.

The foundational stem paper~\cite{GP} provides the appropriate algebraic
framework. The inverse work \cite{CSS,CPSS} supplies local and integral
constructions; \cite{CSS} already contains arctangent primitives for
sphere-integrated Cauchy kernels, and the displayed inverse theorem in
\cite{CPSS} is on a rectangle. The broader spherical-monogenic integral
construction in \cite{CSSSpherical} was also checked at the level of its
published abstract and statement of scope; its full text was not used
as a premise. We do not identify any of these established constructions
as a newly solved inverse-Fueter conjecture.

The singularities in \cite{GPSt} concern slice-regular functions over
alternative $*$-algebras. Our singularity theorem instead concerns axial
\emph{monogenic} fields and includes affine-monodromy modes that cannot
be produced from single-valued slice stems on the punctured base.
The distinction of function class is essential to any comparison of
Laurent theories.

The August 2026 preprint~\cite{DP2026} studies a different inverse problem
for a factor of the Fueter operator, including weak primitives and
axially polyanalytic functions. Its Section 6 was inspected for the
role of domain restrictions. The present paper makes no claim to solve
that different problem; its operator is $\Lap$ on slice stems, and its
principal issue is single-valuedness on the unchanged domain.

The literature search, performed on 21 September 2026, included combinations
of ``inverse Fueter'', ``periods'', ``monodromy'', ``multiply connected'',
``global primitive'', and ``spherical singularities'', together with
backward comparison to the cited inverse papers. It did not locate an
explicit counterpart of the combined period splitting, holomorphic
residue constraint, and $8N$ spherical principal-part classification.
This is a limited, reproducible statement about the search, not a proof
of absence from all published literature. The natural next external
check is specialist comparison with the broader Vekua and generalized
analytic-function literature.

\subsection{What has been proved here}
The global criterion does not rely on a conjectural surjectivity theorem:
it follows from four displayed first-order equations. Every finite
period assignment is realized by explicit fields with verified
monodromy. The singularity normal form follows from that global
criterion on a punctured disk and ordinary Laurent theory. Its
dimension count uses uniqueness of the principal coefficients, not a
formal parameter count alone. The exact asymptotic constants prove the
sharpness and the absence of cancellation at top order.

Accordingly, the unresolved issue for the proposed new results is
external priority and independent checking, not an omitted step in the
stated proofs. Symbolic and numerical checks supplement the proofs;
no formal proof-assistant verification or peer review is claimed.

\subsection{Further mathematical questions}
For infinitely connected bases, the explicit period criterion still
holds, and the range is still closed. The finite sum of logarithmic
modes used here no longer proves that every homomorphism to $\HH^2$ is
realizable. A locally convergent construction with quantitative control
would extend the split theorem beyond finite topology. This is a
question left by this paper, not asserted to be an established open
problem in the literature.

For higher-degree axial monogenics and higher odd-dimensional Fueter
maps, known local kernels are larger polynomial spaces. It is natural
to seek an explicit polynomial-valued monodromy theory, higher moment
periods, and corresponding constrained Laurent coefficients. The
quaternionic case identifies exactly which ingredients such a theory
would require: canonical potentials, a local polynomial kernel, and
single-valued images of logarithms multiplied by kernel polynomials.
No higher-dimensional analogue is asserted without those calculations.

A different question is whether useful boundary-space versions admit
bounded period-corrected inverses with sharp constants. The compact
estimates here avoid boundary regularity assumptions and already give
continuous reconstruction, but they do not settle an optimal Hardy,
Bergman, or Sobolev-space inverse theory.

\appendix
\section{Algebraic consistency checks}
\label{sec:checks}

For the two basic modes, differentiation of
\eqref{eq:V}--\eqref{eq:U} gives identically
\[
 P_x-Q_r-2Q/r=0,\qquad Q_x+P_r=0.
\]
Substitution in \eqref{eq:T} yields
\begin{align*}
 J_{\Vker_\zeta}&=-\frac1{2\pi\iota(z-\zeta)^2},\\
 J_{\Uker_\zeta}&=\frac1{2\pi\iota(z-\zeta)}
                    -\frac\zeta{2\pi\iota(z-\zeta)^2}.
\end{align*}
The first has complex residue zero, yet
$-\int zJ_{\Vker_\zeta}\dd z=1$. The second has
$\int J_{\Uker_\zeta}\dd z=1$, while its first moment integrates to
zero. These calculations check the ordering of the two obstruction
coordinates and the minus sign in the second moment.

The encoding's regular kernel can also be checked directly. For
$R_{e,h}$ in \eqref{eq:T-kernel},
\[
 P_x=0,\qquad Q_x=-2e/r^2,
\]
so $-P+rQ_x=0$ and $J_R=0$. The Vekua equations hold, and its explicit
stem $\iota(ze+h)$ has second derivative zero. This checks that no
complex-affine freedom has been inadvertently included in $\ker\Fu$.

Finally, the critical double-pole coefficient
$j_{-2}=-(\zeta a+b)/(2\pi\iota)$ has
$\ReC j_{-2}=-va/(2\pi)$. Thus
$(\iota/v)\ReC j_{-2}=a/(2\pi\iota)=j_{-1}$, confirming the local
residue constraint with the conventions used throughout.

\begin{thebibliography}{12}
\addcontentsline{toc}{section}{References}

\bibitem{Supplied}
\emph{Classical Complex Theorems over the Quaternions: Slice-regular and
Cauchy--Fueter analogues, with proofs}.
User-supplied expository manuscript, 20 September 2026,
\texttt{quaternionic\_analysis(1).tex}, especially the section
``Fueter's mapping theorem and a constructive local inverse.''

\bibitem{GS}
G.~Gentili and D.~C.~Struppa,
\emph{A new theory of regular functions of a quaternionic variable},
Advances in Mathematics \textbf{216} (2007), 279--301.
\href{https://doi.org/10.1016/j.aim.2007.05.010}{doi:10.1016/j.aim.2007.05.010}.

\bibitem{GP}
R.~Ghiloni and A.~Perotti,
\emph{Slice regular functions on real alternative algebras},
Advances in Mathematics \textbf{226} (2011), 1662--1691.
\href{https://arxiv.org/abs/1008.4318}{arXiv:1008.4318}.

\bibitem{CSS}
F.~Colombo, I.~Sabadini, and F.~Sommen,
\emph{The inverse Fueter mapping theorem},
Communications on Pure and Applied Analysis \textbf{10} (2011), 1165--1181.
\href{https://doi.org/10.3934/cpaa.2011.10.1165}{doi:10.3934/cpaa.2011.10.1165}.

\bibitem{CSSSpherical}
F.~Colombo, I.~Sabadini, and F.~Sommen,
\emph{The inverse Fueter mapping theorem in integral form using spherical
monogenics}, Israel Journal of Mathematics \textbf{194} (2013), 485--505.
\href{https://doi.org/10.1007/s11856-012-0090-4}{doi:10.1007/s11856-012-0090-4}.

\bibitem{CPSS}
F.~Colombo, D.~Pe\~na Pe\~na, I.~Sabadini, and F.~Sommen,
\emph{A new integral formula for the inverse Fueter mapping theorem},
Journal of Mathematical Analysis and Applications \textbf{417} (2014),
112--122.
\href{https://doi.org/10.1016/j.jmaa.2014.03.016}{doi:10.1016/j.jmaa.2014.03.016};
\href{https://arxiv.org/abs/1302.0685}{arXiv:1302.0685}.

\bibitem{GPSt}
R.~Ghiloni, A.~Perotti, and C.~Stoppato,
\emph{Singularities of slice regular functions over real alternative
$*$-algebras}, Advances in Mathematics \textbf{305} (2017), 1085--1130.
\href{https://arxiv.org/abs/1606.03609}{arXiv:1606.03609}.

\bibitem{DP2026}
A.~De Martino and S.~Pinton,
\emph{A variation of the inverse Fueter theorem and the generalized
polyanalytic Cauchy--Kovalevskaya extension of order 2},
preprint (2026), version 1, 7 August 2026.
\href{https://arxiv.org/abs/2608.07711}{arXiv:2608.07711}.

\end{thebibliography}
\end{document}
